\item Considere un rayo de luz que viaja entre el aire y un diamante cortado en la forma mostrada en la figura P34.28.
\begin{enumerate}
    \item Encuentre el ángulo crítico para la reflexión interna total de la luz incidente en el diamante en la interfaz entre el diamante y el aire exterior.
    \item Considere que el rayo de luz incide normalmente en la superficie superior del diamante, como se muestra en la figura. Demuestre que la luz que viaja hacia el punto P en el diamante está totalmente reflejada.
    \item \textbf{¿Qué pasaría si?} Suponga que el diamante se sumerge en agua. ¿Cuál es el ángulo crítico en la interfaz diamante-agua?
    \item Cuando el diamante se sumerge en agua, ¿el rayo de luz que entra por la superficie superior en la figura se somete a una reflexión interna total en P? Explique.
    \item Si el rayo de luz que entra en el diamante sigue siendo vertical, ¿en qué dirección debe rotarse el diamante en el agua alrededor de un eje perpendicular a la página a través de O para que la luz salga del diamante en P?
    \item ¿Con qué ángulo de rotación la luz saldrá primero del diamante en el punto P?
\end{enumerate}